\section*{Erd\H{o}s Problem 690: densities of the \(k\)-th prime factor}

\paragraph{Statement and claim.}
Let \(d_k(p)\) be the density of integers whose \(k\)-th smallest
distinct prime factor is \(p\). The universal unimodality assertion
is false: it fails for \(k=4\). We also prove it for \(k=1,2,3\).
This reconstructs results of Cambie,
\emph{Resolution of Erd\H{o}s' problems about unimodularity},
\url{https://arxiv.org/abs/2501.10333}, with a direct monotonicity
argument for the first three cases.
The existing GPT-5.2 attempt contains the same \(k=4\) witness;
the rational computations are re-derived here.
See \url{https://www.erdosproblems.com/690}.

\paragraph{Exact density formula.}
Divisibility by distinct primes is independent on a complete
residue system modulo their product. Thus the relevant density
exists, and
\[
 d_k(p)=\frac1p[z^{k-1}]
       \prod_{q<p}\left(1-\frac1q+\frac zq\right)
       =\frac{C_p}{p}e_{k-1}(p),
\]
where
\[
 C_p=\prod_{q<p}(1-1/q),\qquad
 \prod_{q<p}\left(1+\frac z{q-1}\right)
   =\sum_{r\geq0}e_r(p)z^r.
\]
All products here are finite. If \(p<q\) are consecutive primes,
then
\[
 C_q=C_p\frac{p-1}{p},\qquad
 e_r(q)=e_r(p)+\frac{e_{r-1}(p)}{p-1}.
\]
Whenever \(e_{k-1}(p)>0\), it follows that
\[
 \frac{d_k(q)}{d_k(p)}
   =\frac{p-1+e_{k-2}(p)/e_{k-1}(p)}q.                         \tag{1}
\]

\paragraph{The first three cases.}
For \(k=1\), directly \(d_1(q)/d_1(p)=(p-1)/q<1\).
For \(k=2\), \(d_2(2)=0\), and for \(p\geq3\) we have
\(e_1(p)=\sum_{\ell<p}1/(\ell-1)\geq1\).
Equation (1) is then at most \(p/q<1\), so the unique peak
occurs at \(p=3\).

For \(k=3\), the initial nonzero term is at \(p=5\).
At that point \(e_1=3/2\), \(e_2=1/2\), so \(3e_2=e_1\).
Adding another variable \(a=1/(\ell-1)>0\) changes these to
\(e'_1=e_1+a\), \(e'_2=e_2+ae_1\), and
\[
 3e'_2-e'_1=(3e_2-e_1)+a(3e_1-1)\geq0.
\]
Hence \(e_1(p)/e_2(p)\leq3\) for every \(p\geq5\).
Consecutive odd primes have \(q\geq p+2\), so (1) is at most
\((p+2)/q\leq1\). The two equal peaks at \(5,7\) have value
\(1/30\); all subsequent terms decrease. This proves unimodality
for \(k=1,2,3\).

\paragraph{A complete counterexample for \(k=4\).}
Expanding the finite products gives
\[
 d_4(13)=\frac{31}{5005},\qquad
 d_4(17)=\frac{206}{36465},\qquad
 d_4(19)=\frac{1308}{230945}.
\]
These can also be computed without fractions until the final step:
the coefficient of \(z^3\) in
\(\prod_{q<p}(q-1+z)\), divided by \(p\prod_{q<p}q\),
is exactly \(d_4(p)\).
The strict comparisons reduce to
\[
 31\cdot36465=1130415>1031030=206\cdot5005,
\]
\[
 206\cdot230945=47574670<47696220=1308\cdot36465.
\]
Thus the sequence decreases from \(13\) to \(17\) and increases
from \(17\) to \(19\), which no unimodal sequence can do.

\paragraph{Scope.}
This completely disproves the assertion for arbitrary fixed \(k\)
and independently proves the positive cases \(1\leq k\leq3\).
It does not purport to classify every individual \(k\geq5\);
the additional computational cases in Cambie's paper are not needed
for the disproof.

\paragraph{Completion estimate.}
\textbf{COMPLETION: 100\%}
